% ej_coord_R2.tex
%
% Copyright (C) 2022--2026 José A. Navarro Ramón <janr.devel@gmail.com>
% Licencia del código GPLv2
% Licencia Creative Commons Recognition Non-Commercial Share-alike.
% (CC-BY-NC-SA)

\chapter{Ejercicios en \mathinhead{\symbb{R}^2}{espr2}}

\section{Ejercicio 1}
\noindent Supongamos un cierto vector en coordenadas rectangulares,
$\vvv{v} = 3\,\uvec{\i} + 4\,\uvec{\j}$.\\
a) Calcula el cuadrado de su módulo.\\
b) Nos proponemos realizar un cambio de base desde el sistema rectangular a otro oblicuo,
con un ángulo entre ejes de \qty{60}{\degree}. Calcula la expresión de los vectores
unitarios de la base nueva. ¿Cuál es su módulo?
Escribe la matriz $\mmm{M}$ de cambio de base y su inversa.\\
c) Obtén el tensor métrico en el sistema oblicuo anterior.\\
d) Halla las componentes contravariantes del vector.\\
e) ¿Cuáles son las componentes covariantes?\\
f) Halla el cuadrado del módulo del vector en el sistema oblicuo.\\

\noindent
\textbf{Solución}

\noindent
a) El cuadrado del módulo del vector es
\[
  v^2 = v_x^2 + v_y^2 = 3^2 + 4^2 = 25
\]

\noindent
b) Utilizando la figura \ref{fig:R2-coord_oblicuas_e1e2}, la nueva base estaría formada
por dos vectores unitarios, que en función del sistema rectangular son
\begin{align*}
  \vvv{e}_1 &= \uvec{\i}
  \longrightarrow |\vvv{e}_1| = 1\\ 
  \vvv{e}_2 &= \cos\pi/3\,\uvec{\i} + \sin\pi/3\,\uvec{\j}
               = \dfrac{1}{2}\,\uvec{\i} + \dfrac{\sqrt{3}}{2}\,\uvec{\j}
               \longrightarrow |\vvv{e}_2| = \cos^2\pi/3 + \sin^2\pi/3 = 1
\end{align*}

La nueva base sería
\[
  \tilde{B} = \set{\vvv{e}_1, \vvv{e}_2}
  = \Set{\uvec{\i}, \dfrac{1}{2}\,\uvec{\i} + \dfrac{\sqrt{3}}{2}\,\uvec{\j}}
  = \Set{(1,0),\, (1/2, \sqrt{3}/2)}
\]

De acuerdo con esto, la matriz de cambio de base estaría formada por los vectores de la
nueva base en función de la antigua, dispuestos en columnas
\[
  \mmm{M}
  = \begin{pmatrix}
    1 & 1/2\\
    0 & \sqrt{3}/2
  \end{pmatrix}
\]

{\footnotesize
Comprobación
\[
  \tilde{B} = B \mmm{M}
\]

Los vectores de la base los disponemos en filas
\[
  \begin{pmatrix}
    \vvv{e}_1 & \vvv{e}_2
  \end{pmatrix}
  =
  \begin{pmatrix}
    \uvec{\i} & \uvec{\j}
  \end{pmatrix}
  \,
  \begin{pmatrix}
    1 & 1/2\\
    0 & \sqrt{3}/2
  \end{pmatrix}
  =
  \begin{pmatrix}
    \uvec{\i} & \dfrac{1}{2}\,\uvec{\i} + \dfrac{\sqrt{3}}{2}\,\uvec{\j}
  \end{pmatrix}
\]
}
La inversa de la matriz de cambio de base se calcula multiplicando la inversa del
determinante de la matriz por la adjunta de su transpuesta. El resultado es
\[
  \mmm{M}^{-1}
  = \dfrac{1}{\det(\mmm{M})}\,\text{adj}\left(\mmm{M}^\transp\right)
  = \dfrac{1}{\sqrt{3}/2}
  \begin{pmatrix}
    \sqrt{3}/2 & -1/2\\[1ex]
    0 & 1
  \end{pmatrix}
  =
  \begin{pmatrix}
    1 & -\sqrt{3}/3\\[1ex]
    0 & 2\sqrt{3}/3
  \end{pmatrix}
\]

\noindent
c) El tensor métrico en el sistema oblicuo, con $\alpha = \pi/3\,\unit{\radian}$ es
\[
  \mmm{g}_{ij}
  =\begin{pmatrix}
    g_{11} & g_{12}\\
    g_{21} & g_{22}
  \end{pmatrix}
  = \begin{pmatrix}
    \vvv{e}_1\cdot\vvv{e}_1 & \vvv{e}_1\cdot\vvv{e}_2\\
    \vvv{e}_2\cdot\vvv{e}_1 & \vvv{e}_2\cdot\vvv{e}_2\\
  \end{pmatrix}
  = \begin{pmatrix}
    1 & \cos\pi/3\\
    \cos\pi/3 & 1
  \end{pmatrix}
  =
  \begin{pmatrix}
    1 & 1/2\\
    1/2 & 1
  \end{pmatrix}
\]

Su inversa es
\[
  \mmm{g}^{ij}
  = \dfrac{1}{\det(\mmm{g}_{ij})}\, \text{adj}(\mmm{g}_{ij}^\transp)
  = \dfrac{1}{3/4}
  \begin{pmatrix}
    1 & -1/2\\
    -1/2 & 1
  \end{pmatrix}
  = \begin{pmatrix}
    4/3 & -2/3\\
    -2/3 & 4/3
  \end{pmatrix}
\]

\noindent
d) Para hallar las componentes contravariantes, debemos aplicar la inversa de la matriz
de cambio de base
\[
  \left[\tilde{x}^i\right]_{\tilde{B}} = \mmm{M}^{-1}\, \left[x^i\right]_{B}
\]
donde las componentes contravarianes en la base rectangular $[x^i]_{B}$ son iguales
que las covariantes y valen $(3, 4)$.

\[
  \begin{pmatrix}
    x^1\\
    x^2
  \end{pmatrix}
  = \begin{pmatrix}
    1 & -\dfrac{\sqrt{3}}{3}\\[1.5ex]
    0 & \dfrac{2\sqrt{3}}{3}
  \end{pmatrix}
  \,
  \begin{pmatrix}
    3\\
    4
  \end{pmatrix}
  = \begin{pmatrix}
    3 - \dfrac{4\sqrt{3}}{3}\\[1.5ex]
    \dfrac{8\sqrt{3}}{3}
  \end{pmatrix}
\]

\noindent
e) Las componentes covariantes se obtienen bajando índices
\[
  x_i = g_{ij}\,x^j
\]
Desarrollando la expresión para $x_1$ y $x_2$ y teniendo en cuenta la métrica,
$g_{11} = 1$, $g_{12}=g_{21}=1/2$ y $g_{22}=1$
\begin{align*}
  x_1
  &= g_{11}x^1 + g_{12}x^2 = x^1 + \dfrac{1}{2}x^2
    = 3 - \cancelout{\dfrac{4\sqrt{3}}{3}} + \cancelout{\dfrac{1}{2}\dfrac{8\sqrt{3}}{3}}
    = 3\\
  x_2
  &= g_{21}x^1 + g_{22} x^2 = \dfrac{1}{2}x^1 + x^2
    = \dfrac{1}{2}\left(3-\dfrac{4\sqrt{3}}{3}\right) + \dfrac{8\sqrt{3}}{3}
    = \dfrac{3}{2} + 2\sqrt{3}
\end{align*}

\noindent
f) El cuadrado del módulo del vector calculado con las componentes oblicuas debe ser
$25$, pues tiene que dar en mismo resultado que en rectangulares (apartado a).
\begin{align*}
  \vvv{v}\cdot\vvv{v}
  &=
    x^ix_i
    = x^1 x_1 + x^2 x_2
    = \left(3-\dfrac{4\sqrt{3}}{3}\right) 3
    + \dfrac{8\sqrt{3}}{3} \left(\dfrac{3}{2} + 2\sqrt{3}\right)\\
  &=
    9 - 4\sqrt{3}
    + \dfrac{\cancelout{3}\cdot 8\sqrt{3}}{2\cdot\cancelout{3}}
    + \dfrac{2\cdot 8\cdot \cancelout{3}}{\cancelout{3}}
    = 9 - \cancelout{4\sqrt{3}} + \cancelout{4\sqrt{3}} + 16
    = 25
\end{align*}

\newpage
\subsection{Ejercicio 2}
\noindent Supongamos un cierto vector en coordenadas rectangulares,
$\vvv{v} = 2\,\uvec{\i} + \sqrt{3}\,\uvec{\j}$.\\
a) Calcula el cuadrado del módulo de $\vvv{v}$.\\
b) Queremos utilizar un sistema de coordenadas oblicuo, cuyos ejes forman
$\qty{60}{\degree}$ entre sí y haciendo coincidir los ejes $x$ y $x'$.
Calcula los vectores de la base $B'=\set{\vvv{e}_1, \vvv{e}_2}$ en el sistema oblicuo
en función de la base del sistema cartesiano ortogonal $B=\set{\uvec{\i},\,\uvec{\j}}$.
Halla sus módulos y comprueba que son vectores unitarios.\\
c) Halla la matriz de cambio de base del sistema cartesiano ortogonal
$B=\set{\uvec{\i},\,\uvec{\j}}$ a la del sistema de coordenadas oblicuas, $B'$.
Calcula también su inversa.\\
d) Obtén las componentes contravariantes del vector $\vvv{v}$ en el sistema de
coordenadas oblicuas, mediante una transformación de componentes.\\
e) Halla ahora las componentes, mediante transformación de los vectores de la base $B$.\\
f) Calcula el producto escalar de los vectores de la base covariante
$B'=\set{\vvv{e}_1, \vvv{e}_2}$ en función de la base en coordenadas ortogonales.\\
g) Obtén el tensor métrico en el sistema oblicuo y su inversa.\\
h) Calcula los vectores de la base contravariante $\set{\vvv{e}^1,\vvv{e}^2}$ en función
de la base rectangular $\set{\uvec{\i},\,\uvec{\j}}$ y halla sus módulos, comprobando
que no son vectores unitarios\\
i) Halla el producto escalar entre los elementos de la base covariante y contravariante.\\
j) Halla las componentes covariantes del vector.\\
k) Halla el cuadrado del módulo del vector en el sistema oblicuo.\\
l) Relaciona el apartado anterior, con el hecho de que los vectores covariantes son
elementos $\omega$ del espacio dual de los vectores $\vvv{v}$ del espacio tangente de $\symbb{R}^2$, y que estas $\omega$ son aplicaciones lineales que al actuar sobre los
vectores del espacio tangente, producen un escalar. Halla el escalar para nuestro vector
$\vvv{v}$.\\

\noindent
\textbf{Solución}

\noindent
a) Cuadrado del módulo de $\vvv{v} = 2\,\uvec{\i} + \sqrt{3}\,\uvec{\j}$
\[
  |\vvv{v}|^2 = v^2 = v_x^2 + v_y^2 = 2^2 + \left(\sqrt{3}\right)^2 = 4 + 3 = 7
\]

\noindent 
b) Matriz de cambio de base de $B=\set{\uvec{\i},\,\uvec{\j}}$ a la de coordenadas
oblicuas que forman $\qty{60}{\degree}$, haciendo coincidir los ejes $x$ y $x'$.
Y también su inversa.

En la figura \ref{fig:R2-coord_oblicuas_e1e2} se puede observar la relación entre estos
vectores, y además podemos comprobar que son unitarios
\begin{align*}
  \vvv{e}_1
  &=
    \uvec{\i}
    \longrightarrow
    |\vvv{e}_1| = |\uvec{\i}|
    = 1\\
  \vvv{e}_2
  &=
    \cos(\pi/3)\,\uvec{\i} + \sin(\pi/3)\,\uvec{\j}
    = \dfrac{1}{2}\,\uvec{\i} + \dfrac{\sqrt{3}}{2}\,\uvec{\j}
    \longrightarrow
    |\vvv{e}_2| = \left|\dfrac{1}{2}\,\uvec{\i}+\dfrac{\sqrt{3}}{2}\,\uvec{\j}\right|
    = \dfrac{1}{4} + \dfrac{3}{4}
    = 1
\end{align*}

\noindent
c) Matriz de cambio de base $B$ a $B'$ y su inversa.

La transformación de vectores de la base es $B'=B \mmm{M}$.
Del apartado anterior se puede comprobar que
\[
  \begin{pmatrix}
    \vvv{e}_1 & \vvv{e}_2
  \end{pmatrix}
  = \begin{pmatrix}
    \uvec{\i} & \uvec{\j}
  \end{pmatrix}
  \,\begin{pmatrix}
    1 & 1/2\\
    0 & \sqrt{3}/2
  \end{pmatrix}
\]

Luego, la matriz de cambio de base de $B$ a $B'$ es
\[
  \mmm{M}
  = \begin{pmatrix}
    1 & 1/2\\
    0 & \sqrt{3}/2
  \end{pmatrix}
\]

Su inversa es
\[
  \mmm{M}^{-1}
  = \dfrac{1}{\det(\mmm{M})}\,
  \text{adj}(\mmm{M}^\transp)
  = \dfrac{1}{\sqrt{3}/2}
  \begin{pmatrix}
    \sqrt{3}/2 & -1/2\\
    0 & 1
  \end{pmatrix}
  = \begin{pmatrix}
    1 & -\sqrt{3}/3\\
    0 & 2\sqrt{3}/3
  \end{pmatrix}
\]

\noindent
d) Componentes contravariantes de $\vvv{v}$.

Estas componentes se calculan con la inversa de la matriz de cambio de base
\[
  \begin{pmatrix}
    x^1\\
    x^2
  \end{pmatrix}
  = \begin{pmatrix}
    1 & -\sqrt{3}/3\\
    0 & 2\sqrt{3}/3
  \end{pmatrix}
  \,\begin{pmatrix}
    x\\
    y
  \end{pmatrix}
\]
\[
  \begin{pmatrix}
    x^1\\
    x^2
  \end{pmatrix}
  = \begin{pmatrix}
    1 & -\sqrt{3}/3\\
    0 & 2\sqrt{3}/3
  \end{pmatrix}
  \,\begin{pmatrix}
    2\\[0.4ex]
    \sqrt{3}
  \end{pmatrix}
  = \begin{pmatrix}
    1\\
    2
  \end{pmatrix}
\]

Las coordenadas contravariantes en el nuevo sistema son
\begin{align*}
  x^1 = 1\\
  x^2 = 2
\end{align*}

\noindent
e) Componentes contravariantes mediante transformación de vectores de la base $B$.

Necesitamos obtener los vectores de la base antigua $B=\set{\uvec{\i},\uvec{\j}}$ en
función de la nueva $B'=\set{\vvv{e}_1,\vvv{e}_2}$. Para ello utilizaremos la inversa de
la matriz de cambio de base
\[
  \begin{pmatrix}
    \uvec{\i} & \uvec{\j}
  \end{pmatrix}
  =
  \begin{pmatrix}
    \vvv{e}_1 & \vvv{e}_2
  \end{pmatrix}
  \begin{pmatrix}
    1 & -\dfrac{\sqrt{3}}{3}\\[1.4ex]
    0 & \dfrac{2\sqrt{3}}{3}
  \end{pmatrix}
  =
  \begin{pmatrix}
    \vvv{e}_1 & -\dfrac{\sqrt{3}}{3}\vvv{e}_1+\dfrac{2\sqrt{3}}{3}\vvv{e}_2
  \end{pmatrix}
\]

La transformación de vectores resulta
\begin{align*}
  \uvec{\i} &= \vvv{e}_1\\
  \uvec{\j} &= -\dfrac{\sqrt{3}}{3}\vvv{e}_1 + \dfrac{2\sqrt{3}}{3}\vvv{e}_2
\end{align*}

Por último, se sustituyen estos vectores de $B$ en la expresión de $\vvv{v}$
\[
  \vvv{v}
  =
  2\,\uvec{\i} + \sqrt{3}\,\uvec{\j}
  = 2 \vvv{e}_1
  +\sqrt{3}\left(-\dfrac{\sqrt{3}}{3}\vvv{e}_1+\dfrac{2\sqrt{3}}{3}\vvv{e}_2\right)
  =
  2\vvv{e}_1 - \vvv{e}_1 + 2\vvv{e}_2
  = \vvv{e}_1 + 2\vvv{e}_2 
\]

El vector expresado con componentes contravariantes es
\[
  \vvv{v} = x^i\vvv{e}_i = x^1\vvv{e}_1 + x^2\vvv{e}_2
\]

Comparando con nuestro resultado
\[
  \vvv{v} = x^1\vvv{e}_1 + x^2\vvv{e}_2 = 1\vvv{e}_1 + 2\vvv{e}_2
\]

Estas componentes son, obviamente, las mismas que en el apartado anterior
\begin{align*}
  x^1 &= 1\\
  x^2 &= 2
\end{align*}

\noindent
f) Producto escalar de los vectores de la base covariante $B=\set{\vvv{e}_1,\vvv{e}_2}$.
\begin{align*}
  \vvv{e}_1\cdot\vvv{e}_1
  &=
    \uvec{\i}\cdot\uvec{\i}
    = 1\\
  \vvv{e}_1\cdot\vvv{e}_2
  &=
    \vvv{e}_2\cdot\vvv{e}_1
    = \uvec{\i}\cdot\left(\dfrac{1}{2}\,\uvec{\i}+\dfrac{\sqrt{3}}{2}\,\uvec{\j}\right)
    = \dfrac{1}{2}\\
  \vvv{e}_2\cdot\vvv{e}_2
  &=
    \left(\dfrac{1}{2}\,\uvec{\i}+\dfrac{\sqrt{3}}{2}\,\uvec{\j}\right)
    \cdot
    \left(\dfrac{1}{2}\,\uvec{\i}+\dfrac{\sqrt{3}}{2}\,\uvec{\j}\right)
    = \left(\dfrac{1}{2}\right)^2 + \left(\dfrac{\sqrt{3}}{2}\right)^2
    = \dfrac{1}{4} + \dfrac{3}{4}
    = 1
\end{align*}

\noindent
g) Tensor métrico y su inverso.

El tensor métrico sería
\[
  \mmm{g}_{ij}
  =
  \begin{pmatrix}
    g_{11} & g_{12}\\
    g_{21} & g_{22}
  \end{pmatrix}
  =
  \begin{pmatrix}
    \xhat{e}_1\cdot\xhat{e}_1 & \xhat{e}_1\cdot\xhat{e}_2\\
    \xhat{e}_2\cdot\xhat{e}_1 & \xhat{e}_2\cdot\xhat{e}_2
  \end{pmatrix}
  =
  \begin{pmatrix}
    1 & 1/2\\
    1/2 & 1
  \end{pmatrix}
\]

La inversa del tensor métrico es
\[
  \mmm{g}^{ij}
  = \dfrac{1}{\det(\mmm{g}_{ij})}\, \text{adj}(\mmm{g}_{ij}^\transp)
  = \dfrac{1}{3/4}
  \begin{pmatrix}
    1 & -1/2\\
    -1/2 & 1
  \end{pmatrix}
  = \begin{pmatrix}
    4/3 & -2/3\\
    -2/3 & 4/3
  \end{pmatrix}
\]

\noindent
h) Base contravariante en función de $\uvec{\i}$ y $\uvec{\j}$.

En el apartado d) se calcularon los vectores de la base covariante, resultando
\begin{align*}
  \vvv{e}_1 &= \uvec{\i}\\
  \vvv{e}_2 &= \dfrac{1}{2}\,\uvec{\i} + \dfrac{\sqrt{3}}{2}\,\uvec{\j}
\end{align*}

Subiendo índices obtendríamos los vectores de la base contravariante
\[
  \vvv{e}^i = g^{ij}\vvv{e}_j
\]
\begin{align*}
  \vvv{e}^1
  &=
    g^{11}\vvv{e}_1 + g^{12}\vvv{e}_2
    = \dfrac{4}{3}\vvv{e}_1 - \dfrac{2}{3}\vvv{e}_2
    = \dfrac{4}{3}\,\uvec{\i}
    - \dfrac{2}{3}\,\left(\dfrac{1}{2}\,\uvec{\i}+\dfrac{\sqrt{3}}{2}\,\uvec{\j}\right)
    = \uvec{\i} - \dfrac{\sqrt{3}}{3}\,\uvec{\j}\\
  \vvv{e}^2
  &=
    g^{21}\vvv{e}_1 + g^{22}\vvv{e}_2
    = -\dfrac{2}{3}\vvv{e}_1 + \dfrac{4}{3}\vvv{e}_2
    = -\dfrac{2}{3}\,\uvec{\i}
    + \dfrac{4}{3}\,\left(\dfrac{1}{2}\,\uvec{\i}+\dfrac{\sqrt{3}}{2}\,\uvec{\j}\right)
    = \dfrac{2\sqrt{3}}{3}\,\uvec{\j}  
\end{align*}

Los vectores de la base contravariante son
\begin{align*}
  \vvv{e}^1 &= \uvec{\i} - \dfrac{\sqrt{3}}{3}\,\uvec{\j}\\
  \vvv{e}^2 &= \dfrac{2\sqrt{3}}{3}\,\uvec{\j}
\end{align*}

Calculamos sus módulos y comprobamos que no son vectores unitarios
\begin{align*}
  \left|\vvv{e}^1\right|
  &=
    1^2 + \left(\dfrac{\sqrt{3}}{3}\right)^2
    = 1 + \dfrac{1}{3}
    = \dfrac{4}{3}\\
  \left|\vvv{e}^2\right|
  &=
    \left(\dfrac{2\sqrt{3}}{3}\right)^2
    = \dfrac{4}{3}
\end{align*}

\noindent
i) Producto escalar entre elementos de las bases covariante y contravariante
\begin{align*}
  \vvv{e}^1\cdot\vvv{e}_1
  &=
    \left(\uvec{\i} - \dfrac{\sqrt{3}}{3}\,\uvec{\j}\right)
    \cdot
    \uvec{\i}
    = 1\\
  \vvv{e}^2\cdot\vvv{e}_2
  &=
    \left(\dfrac{2\sqrt{3}}{3}\,\uvec{\j}\right)
    \cdot
    \left(\dfrac{1}{2}\,\uvec{\i} + \dfrac{\sqrt{3}}{2}\,\uvec{\j}\right)
    = 1
\end{align*}
Y los productos cruzados
\begin{align*}
  \vvv{e}^1\cdot\vvv{e}_2
  &=
    \left(\uvec{\i} - \dfrac{\sqrt{3}}{3}\,\uvec{\j}\right)
    \cdot
    \left(\dfrac{1}{2}\,\uvec{\i} + \dfrac{\sqrt{3}}{2}\,\uvec{\j}\right)
    = \dfrac{1}{2} - \dfrac{\sqrt{3}}{3}\cdot\dfrac{\sqrt{3}}{2}
    = \dfrac{1}{2} - \dfrac{1}{2}
    = 0\\
  \vvv{e}^2\cdot\vvv{e}_1
  &=
    \left(
    \dfrac{2\sqrt{3}}{3}\,\uvec{\j}\right) \cdot \uvec{\i}
    = 0
\end{align*}

\noindent
j) Componentes covariantes de $\vvv{v}$.

Las componentes covariantes del vector se pueden calcular bajando el índice de las
contravariantes
\[
  x_i = g_{ij} x^j
\]
\begin{align*}
  x_1
  &=
    g_{11} x^1 + g_{12} x^2
    = 1\cdot 1 + \dfrac{1}{2} \cdot 2
    =  2\\
  x_2
  &=
    g_{21} x^1 + g_{22} x^2
    = \dfrac{1}{2} \cdot 1 + 1\cdot 2
    = \dfrac{5}{2}
\end{align*}

\noindent
k) Cuadrado de $\vvv{v}$ en el sistema oblicuo.

Se comprueba que el vector no ha cambiado. Solo lo han hecho sus componentes en el nuevo
sistema de coordenadas
\[
  |\vvv{v}|^2
  = v^2
  = x_ix^i
  = x_1x^1 + x_2x^2
  = 2\cdot 1 + \dfrac{5}{2}\cdot 2
  = 2 + 5
  = 7
\]

\noindent
l) Explica el apartado anterior, con el hecho de que los vectores covariantes son
elementos $\omega$ del espacio dual de los vectores $\vvv{v}$ del espacio tangente de
$\symbb{R}^2$ y que estas $\omega$ son aplicaciones lineales que al actuar sobre un
vector, producen un escalar.

El vector $\vvv{v}$ es un elemento del espacio tangente de $\symbb{R}^2$
\[
  \vvv{v} \in T_p\symbb{R}^2
\]

Podemos establecer una aplicación lineal $\omega$ que, al actuar sobre un vector, produce
un escalar
\[
  \omega: T_p\symbb{R}^2 \rightarrow \symbb{R}
\]

Diremos que $\omega$ es un elemento del \emph{espacio dual} del vector
\[
  \omega \in \left(T_p\symbb{R}^2\right)^*
\]

Los vectores covariantes son estas aplicaciones lineales $\omega = x_ie^i$ y, como hemos
dicho, son elementos del espacio dual, que no es el espacio tangente donde \emph{viven}
los vectores. La métrica permite relacionar los dos espacios (subiendo o bajando índices).

Ahora nos podemos preguntar que si
$\omega = x_1\vvv{e}^1 + x_2\vvv{e}^2 = 2\vvv{e}^1 + \frac{5}{2}\vvv{e}^2$, es una de
tales aplicaciones lineales, debería \emph{comer} un vector como $\vvv{v}$ y producir un
escalar (en este caso un número real).

Veamos en notación de Einstein, que el número real es el cuadrado del módulo del vector
\[
  \omega(\vvv{v}) = x_ie^i\cdot x^ie_i = x_ix^i = v^2
\]

En nuestro ejemplo, tendríamos
\[
  \omega(\vvv{v})
  = \left(2\vvv{e}^1+\dfrac{5}{2}\vvv{e}^2\right)\cdot\left(1\vvv{e}_1+2\vvv{e}_2\right)
  = 2\cdot 1 + \dfrac{5}{2}\cdot 2
  = 2 + 5
  = 7
\]
donde se han tenido en cuenta los productos escalares del apartado anterior.

Recalcamos el hecho de que la aplicación lineal $\omega\in (T_p\symbb{R}^2)^*$
(elemento del espacio dual de $\vvv{v} = 2\,\uvec{\i}+\sqrt{3}\,\uvec{\j}$), que es
$\omega=2\vvv{e}^1+\frac{5}{2}\vvv{e}^2$, cuando se aplica al vector $\vvv{v}$ nos da un
escalar, que resulta ser el cuadrado del módulo del vector


\newpage
\subsection{Ejercicio 3}
Un campo escalar está definido en coordenadas cartesianas ortogonales en $\symbb{R}^2$
por
\[
  \phi(x,y) = xy+2x
\]
Nos interesa utilizar un sistema de coordenadas cartesianas oblicuas, cuyos ejes forman
\qty{60}{\degree} entre sí, donde el eje de abcisas de ambos sistemas coinciden.\\
Calcula lo que se pide, primero en general para cualquier sistema oblicuo, y después
para nuestro sistema oblicuo concreto. Por último, particulariza para el vector
$\vvv{v} = (x,y) = (3,2)$, en los casos en los que sea pertinente.\\
a) El gradiente en cartesianas ortogonales y su cuadrado.\\
%b) La matriz de cambio de base $\mmm{M}$ y su inversa.\\
%c) Las componentes contravariantes en función de las ortogonales.\\
%d) Las componentes ortogonales en función de las contravariantes.\\
%e) Halla la base covariante en función de la base ortogonal y viceversa.\\
%f) Obtén la métrica en coordenadas oblicuas y su inversa.\\
%g) Calcula las componentes covariantes del gradiente.\\
%h) Halla las componentes contravariantes del gradiente.\\
%i) Calcula el cuadrado del gradiente en coordenadas oblicuas.\\

\noindent
\textbf{Solución}

\noindent
a) Gradiente en cartesianas ortogonales y su cuadrado.

Derivadas parciales
\[
  \dfrac{\partial\phi}{\partial x} = y + 2
  \,;\hspace{1em}
  \dfrac{\partial\phi}{\partial y} = x
\]

Gradiente y su cuadrado
\begin{align*}
  &\vvv{\nabla}\phi
    = \dfrac{\partial\phi}{\partial x}\,\uvec{\i}
    + \dfrac{\partial\phi}{\partial y}\,\uvec{\j}
    = (y+2)\,\uvec{\i} + x\,\uvec{\j}\\
  &|\vvv{\nabla}\phi|^2
    = (y+2)^2 + x^2
\end{align*}
 
Gradiente y su cuadrado en el punto $x=3$ e $y=2$
\begin{align*}
  &\vvv{\nabla}\phi(3,2)
    = (y+2)\,\uvec{\i} + x\,\uvec{\j}
    = (2+2)\,\uvec{\i} + 3\,\uvec{\j}
    = 4\,\uvec{\i} + 3\,\uvec{\j}\\
  & |\vvv{\nabla}\phi|^2
    = 4^2 + 3^2
    = 25
\end{align*}

\noindent
%b) Matriz de cambio de base y su inversa.
%
%Mostramos el cambio de base, según la figura \ref{fig:R2-coord_oblicuas_e1e2}, y
%su expresión matricial
%\[
%  \begin{cases}
%    \begin{array}{l}
%      \xhat{e}_1 = \uvec{\i}\\
%      \xhat{e}_2 = \cos\alpha\,\uvec{\i} + \sin\alpha\,\uvec{\j}
%    \end{array}
%  \end{cases}
%  \longrightarrow\hspace{1em}
%  \begin{pmatrix}
%    \xhat{e}_1 & \xhat{e}_2
%  \end{pmatrix}
%  = \begin{pmatrix}
%    \uvec{\i} & \uvec{\j}
%  \end{pmatrix}
%  \,\begin{pmatrix}
%    1 & \cos\alpha\\
%    0 & \sin\alpha
%  \end{pmatrix}
%\]
%
%De aquí, se deduce que la matriz de cambio de base es
%\[
%  \mmm{M}
%  =\begin{pmatrix}
%    1 & \cos\alpha\\
%    0 & \sin\alpha
%  \end{pmatrix}
%\]
%
%La inversa resulta ser
%\begin{align*}
%  \mmm{M}^{-1}
%  &=
%    \dfrac{1}{\det(\mmm{M})}
%    \,\text{adj}(\mmm{M}^\transp)
%    = \dfrac{1}{\sin\alpha}
%    \begin{pmatrix}
%      \sin\alpha & -\cos\alpha\\
%      0 & 1
%    \end{pmatrix}
%    = \begin{pmatrix}
%      1 & -\cot\alpha\\
%      0 & \phantom{-}\csc\alpha
%    \end{pmatrix}
%\end{align*}
%
%En nuestro sistema oblicuo, con $\alpha = \qty{60}{\degree}$
%\begin{align*}
%  \mmm{M}
%  &=
%    \begin{pmatrix}
%      1 & \cos\pi/3\\
%      0 & \sin\pi/3
%    \end{pmatrix}
%    = \begin{pmatrix}
%      1 & 1/2\\
%      0 & \sqrt{3}/2
%    \end{pmatrix}\\
%  \mmm{M}^{-1}
%  &=
%    \begin{pmatrix}
%      1 & -\cot\pi/3\\
%      0 & \csc\pi/3
%    \end{pmatrix}
%    = \begin{pmatrix}
%      1 & -\sqrt{3}/3\\
%      0 & 2\sqrt{3}/3
%    \end{pmatrix}
%\end{align*}
%
%Resumen:
%\[
%  \text{Matriz de cambio de base general e inversa }
%  \begin{cases}
%    \begin{array}{l}
%      \mmm{M}
%      = \begin{pmatrix}
%        1 & \cos\alpha\\
%        0 & \sin\alpha
%      \end{pmatrix}\\
%      \mmm{M}^{-1}
%      = \begin{pmatrix}
%        1 & -\cot\alpha\\
%        0 & \csc\alpha
%      \end{pmatrix}      
%    \end{array}
%  \end{cases}
%\]
%
%\[
%  \text{Matriz de cambio de base $\alpha=\qty{60}{\degree}$ e inversa }
%  \begin{cases}
%    \begin{array}{l}
%      \mmm{M}
%      = \begin{pmatrix}
%        1 & 1/2\\
%        0 & \sqrt{3}/2
%      \end{pmatrix}\\
%      \mmm{M}^{-1}
%      = \begin{pmatrix}
%        1 & -\sqrt{3}/3\\
%        0 & 2\sqrt{3}/3
%      \end{pmatrix}      
%    \end{array}
%  \end{cases}
%\]
%
%\noindent
%c) Componentes contravariantes en función de las ortogonales.
%
%Valores de $x^1$ y $x^2$ para el vector $\vvv{v} = (x,y)$
%\[
%  \begin{pmatrix}
%    x^1\\
%    x^2
%  \end{pmatrix}
%  =
%  \begin{pmatrix}
%    1 & -\cot\alpha\\
%    0 & \csc\alpha
%  \end{pmatrix}
%  \,\begin{pmatrix}
%    x\\
%    y
%  \end{pmatrix}
%  \longrightarrow
%  \begin{cases}
%    \begin{array}{l}
%      x^1 = x-y\cot\alpha\\
%      x^2 = y\csc\alpha
%    \end{array}
%  \end{cases}
%\]
%\[
%  \begin{array}{ll}
%  \text{Para nuestro sistema oblicuo y } \vvv{v} = (x,y)&
%  \begin{cases}
%    \begin{array}{l}
%      x^1 = x-\dfrac{\sqrt{3}}{3}\,y\\
%      x^2 = \dfrac{2\sqrt{3}}{3}\,y
%    \end{array}
%  \end{cases}\\
%  \text{Para nuestro sistema oblicuo y } \vvv{v} = (3,2)&
%  \begin{cases}
%    \begin{array}{l}
%      x^1 = 3-\dfrac{2\sqrt{3}}{3}\\
%      x^2 = \dfrac{4\sqrt{3}}{3}
%    \end{array}
%  \end{cases}
%  \end{array}
%\]
%
%\noindent
%d) Las componentes ortogonales en función de las contravariantes.\\
%\[
%  \begin{pmatrix}
%    x\\
%    y
%  \end{pmatrix}
%  =
%  \begin{pmatrix}
%    1 & \cos\alpha\\
%    0 & \sin\alpha
%  \end{pmatrix}
%  \,\begin{pmatrix}
%    x^1\\
%    x^2
%  \end{pmatrix}
%  \longrightarrow
%  \begin{cases}
%    \begin{array}{l}
%      x = x^1 + x^2\cos\alpha\\
%      y = x^2\sin\alpha
%    \end{array}
%  \end{cases}
%\]
%
%\[
%  \text{En nuestro sistema oblicuo \qty{60}{\degree} }
%  \begin{cases}
%    \begin{array}{l}
%      x = x^1 + \dfrac{1}{2} x^2\\
%      y = \dfrac{\sqrt{3}}{2} x^2
%    \end{array}
%  \end{cases}
%\]
%
%\noindent
%e) Transformación de la base ortogonal a la covariante del sistema oblicuo y su inversa.
%
%Con la matriz de cambio de base $\mmm{M}$, se obtienen los vectores de la base
%covariante
%\[
%  \begin{pmatrix}
%    \xhat{e}_1 & \xhat{e}_2
%  \end{pmatrix}
%  =\begin{pmatrix}
%    \uvec{\i} & \uvec{\j}
%  \end{pmatrix}
%  \,\begin{pmatrix}
%    1 & \cos\alpha\\
%    0 & \sin\alpha
%  \end{pmatrix}
%  \longrightarrow
%  \begin{cases}
%    \begin{array}{l}
%      \xhat{e}_1 = \uvec{\i}\\
%      \xhat{e}_2 = \cos\alpha\,\uvec{\i} + \sin\alpha\,\uvec{\j}
%    \end{array}
%  \end{cases}
%\]
%Estos vectores miden las componentes contravariantes de un vector en $\symbb{R}^2$ en
%coordenadas oblicuas.
%
%La transformación inversa es
%\[
%  \begin{pmatrix}
%    \uvec{\i} & \uvec{\j}                 
%  \end{pmatrix}
%  =\begin{pmatrix}
%    \xhat{e}_1 & \xhat{e}_2
%  \end{pmatrix}
%  \,\begin{pmatrix}
%    1 & -\cot\alpha\\
%    0 & \csc\alpha
%  \end{pmatrix}
%  \longrightarrow
%  \begin{cases}
%    \begin{array}{l}
%      \uvec{\i} = \xhat{e}_1\\
%      \uvec{\j} = -\cot\alpha\,\xhat{e}_1 + \csc\alpha\,\xhat{e}_2
%    \end{array}
%  \end{cases}
%\]
%Resumen: En nuestro sistema oblicuo $\alpha=\qty{60}{\degree}$, tenemos
%\[
%  \begin{array}{ll}
%  \text{Transformación de base ortogonal a oblicua }&
%  \begin{cases}
%    \begin{array}{l}
%      \xhat{e}_1 = \uvec{\i}\\
%      \xhat{e}_2 = \dfrac{1}{2}\,\uvec{\i} + \dfrac{\sqrt{3}}{2}\,\uvec{\j}
%    \end{array}
%  \end{cases}\\[4ex]
%  \text{Transformación de base covariante a ortogonal }&
%  \begin{cases}
%    \begin{array}{l}
%      \uvec{\i} = \xhat{e}_1\\
%      \uvec{\j} = -\dfrac{\sqrt{3}}{3}\,\xhat{e}_1 + \dfrac{2\sqrt{3}}{3}\,\xhat{e}_2
%    \end{array}
%  \end{cases}
%  \end{array}
%\]
%
%\noindent
%f) Métrica y su inversa.
%
%La métrica se calcula mediante el producto escalar de los vectores de la base covariante
%\begin{align*}
%  \mmm{g}_{ij}
%  &=
%    \begin{pmatrix}
%      g_{11} & g_{12}\\
%      g_{21} & g_{22}
%    \end{pmatrix}
%    = \begin{pmatrix}
%      \xhat{e}_1\cdot \xhat{e}_1 & \xhat{e}_1\cdot\xhat{e}_2\\
%      \xhat{e}_2\cdot \xhat{e}_1 & \xhat{e}_2\cdot\xhat{e}_2
%    \end{pmatrix}
%    = \begin{pmatrix}
%      1 & \cos\alpha\\
%      \cos\alpha & 1
%    \end{pmatrix}
%\end{align*}
%
%La métrica inversa es
%\begin{align*}
%  \mmm{g}^{ij}
%  &=
%    \dfrac{1}{\det(g_{ij})}\,
%    \text{adj}\!\left(g_{ij}^\transp\right)
%    = \dfrac{1}{\sin^2\alpha}
%    \begin{pmatrix}
%      1 & -\cos\alpha\\
%      -\cos\alpha & 1
%    \end{pmatrix}
%\end{align*}
%
%\noindent
%Resumen:
%\[
%  \hspace*{-5.5em}
%  \text{Métrica general y su inversa }
%  \begin{cases}
%    \begin{array}{l}
%      \mmm{g}_{ij}
%      = \begin{pmatrix}
%        1 & \cos\alpha\\
%        \cos\alpha & 1
%      \end{pmatrix}\\
%      \mmm{g}^{ij}
%      = \dfrac{1}{\sin^2\alpha}
%      \begin{pmatrix}
%        1 & -\cos\alpha\\
%        -\cos\alpha & 1
%      \end{pmatrix}
%    \end{array}
%  \end{cases}
%\]
%\[
%  \text{Métrica y su inversa para \qty{60}{\degree} }
%  \begin{cases}
%    \begin{array}{l}
%      \mmm{g}_{ij}
%      = \begin{pmatrix}
%        1 & 1/2\\
%        1/2 & 1
%      \end{pmatrix}\\
%      \mmm{g}^{ij}
%      =
%      \dfrac{1}{3/4}
%      \begin{pmatrix}
%        1 & -1/2\\
%        -1/2 & 1
%      \end{pmatrix}
%      = \begin{pmatrix}
%        4/3 & -2/3\\
%        -2/3 & 4/3
%      \end{pmatrix}
%    \end{array}
%  \end{cases}
%\]
%
%\noindent
%g) Componentes covariantes del gradiente.
%
%Del apartado \emph{d)}, sabemos las coordenadas ortogonales a partir de las
%contravariantes. Las sustituimos en las derivadas parciales en cartesianas ortogonales
%de \emph{a)}
%\[
%  \begin{cases}
%    \begin{array}{l}
%      x = x^1 + x^2\,\cos\alpha\\
%      y = x^2\sin\alpha
%    \end{array}
%  \end{cases}
%\]
%
%Expresamos las componentes ortogonales en función de las contravariantes
%\begin{align*}
%  \dfrac{\partial\phi}{\partial x}
%  &=
%    y + 2
%    = x^2\sin\alpha + 2\\
%  \dfrac{\partial\phi}{\partial y}
%  &=
%    x
%    = x^1 + x^2\cos\alpha
%\end{align*}
%
%Calculamos las componentes covariantes del gradiente, teniendo en cuenta las derivadas
%parciales de $x$ e $y$ respecto de $x^1$ y $x^2$, que dejamos como ejercicio
%\begin{align*}
%  \hspace{-14em}
%  \nabla\phi_1
%  &=
%    \dfrac{\partial\phi}{\partial x^1}
%    = \dfrac{\partial\phi}{\partial x}\dfrac{\partial x}{\partial x^1}
%    + \dfrac{\partial\phi}{\partial y}\dfrac{\partial y}{\partial x^1}
%    = (x^2\sin\alpha+2)\cdot 1 + (x^1 + x^2\cos\alpha)\cdot 0\\
%  &=
%    x^2\sin\alpha + 2
%\end{align*}
%\begin{align*}
%  \hspace{1.3em}
%  \nabla\phi_2
%  &=
%    \dfrac{\partial\phi}{\partial x^2}
%    = \dfrac{\partial\phi}{\partial x}\dfrac{\partial x}{\partial x^2}
%    + \dfrac{\partial\phi}{\partial y}\dfrac{\partial y}{\partial x^2}
%    = (x^2\sin\alpha+2)\cdot\cos\alpha + (x^1 + x^2\cos\alpha)\cdot\sin\alpha\\
%  &=
%    x^1\sin\alpha + 2x^2\sin\alpha\cos\alpha + 2\cos\alpha
%\end{align*}
%
%En nuestro sistema oblicuo \qty{60}{\degree}
%\begin{align*}
%  \nabla\phi_1
%  &=
%    x^2\sin\pi/3 + 2
%    = \dfrac{\sqrt{3}}{2} x^2 + 2\\
%  \nabla\phi_2
%  &=
%    x^1\sin\pi/3 + x^2\,2\sin\pi/3\cos\pi/3 + 2\cos\pi/3
%    = \dfrac{\sqrt{3}}{2}x^1 + \dfrac{\sqrt{3}}{2} x^2 + 1
%\end{align*}
%
%Para el sistema oblicuo y el vector $\vvv{v} = (3,2)$
%\begin{align*}
%  \hspace{-7em}
%  \nabla\phi_1
%  &=
%    \dfrac{\sqrt{3}}{2}\cdot \dfrac{4\sqrt{3}}{3} + 2
%    = 2 + 2
%    = 4
%\end{align*}
%\begin{align*}
%  \nabla\phi_2
%  &=
%    \dfrac{\sqrt{3}}{2}\cdot \left(3-2\dfrac{\sqrt{3}}{3}\right)
%    + \dfrac{\sqrt{3}}{2}\cdot\dfrac{4\sqrt{3}}{3} + 1
%    = \dfrac{3\sqrt{3}}{2} - 1 + 2 + 1\\
%  &=
%    2+\dfrac{3\sqrt{3}}{2}
%\end{align*}
%
%\noindent
%Resumen:
%\[
%  \hspace{1em}
%  \text{Sistema oblicuo general }
%  \begin{cases}
%    \begin{array}{l}
%      \nabla\phi_1 = x^2\sin\alpha + 2\\
%      \nabla\phi_2 = x^1\sin\alpha + 2x^2\sin\alpha\cos\alpha + 2 \cos\alpha
%    \end{array}
%  \end{cases}
%\]
%\[
%  \begin{array}{ll}
%    \text{Sistema oblicuo \qty{60}{\degree} } &
%    \begin{cases}
%      \begin{array}{l}
%        \nabla\phi_1 = \dfrac{\sqrt{3}}{2} x^2 + 2\\[2ex]
%        \nabla\phi_2 =\dfrac{\sqrt{3}}{2} x^1 + \dfrac{\sqrt{3}}{2} x^2 + 1
%      \end{array}
%    \end{cases}\\[2ex]
%    \text{Sistema oblicuo y vector } \vvv{v} = (3,2) &
%    \begin{cases}
%      \begin{array}{l}
%        \nabla\phi_1 = 4\\
%        \nabla\phi_2 = 2 + \dfrac{3\sqrt{3}}{2}
%      \end{array}
%    \end{cases}
%  \end{array}
%\]
%
%\noindent
%h) Componentes contravariantes del gradiente.
%
%\vspace*{1ex}
%Se calculan mediante $\nabla\phi^i = g^{ij} \nabla\phi_j$
%\begin{align*}
%  \nabla\phi^1
%  &=
%    g^{11} \nabla\phi_1 + g^{12}\nabla\phi_2\\
%  &=
%    \dfrac{1}{\sin^2\alpha} \left(x^2\,\sin\alpha + 2\right)
%    - \dfrac{\cos\alpha}{\sin^2\alpha}
%    \,\left(x^1\sin\alpha
%    + 2x^2\sin\alpha\cos\alpha
%    + 2\cos\alpha\right)\\
%  &=
%    - x^1\,\cot\alpha
%    + x^2\,\dfrac{1-2\cos^2\alpha}{\sin\alpha}
%    + 2
%\end{align*}
%\begin{align*}
%  \nabla\phi^2
%  &=
%    g^{21} \nabla\phi_1 + g^{22}\nabla\phi_2\\
%  &=
%    -\dfrac{\cos\alpha}{\sin^2\alpha} \left(x^2\,\sin\alpha + 2\right)
%    + \dfrac{1}{\sin^2\alpha}
%    \,\left(x^1\sin\alpha
%    + x^2\,2\sin\alpha\cos\alpha
%    + 2\cos\alpha\right)\\
%  &=
%    x^1\csc\alpha
%    + x^2\left(%
%    2\,\cot\alpha
%    -\cot\alpha
%    \right)
%    - \cancelout{\dfrac{2\cos\alpha}{\sin^2\alpha}}
%    + \cancelout{\dfrac{2\cos\alpha}{\sin^2\alpha}}
%    = x^1\csc\alpha
%    + x^2\cot\alpha
%\end{align*}
%
%Para el sistema oblicuo con $\alpha= \qty{60}{\degree}$
%\begin{align*}
%  \nabla\phi^1
%  &=
%    -x^1 \cot{\pi/3}
%    + x^2\,\dfrac{1-2\cos^2\pi/3}{\sin\pi/3}
%    + 2\\
%  &=
%    -\dfrac{\sqrt{3}}{3}x^1
%    +  \dfrac{1-2\,(1/4)}{\sqrt{3}/2} x^2
%    + \dfrac{2\cos\pi/3}{\sin^2\pi/3}
%    = -\dfrac{\sqrt{3}}{3}x^1
%    + \dfrac{\sqrt{3}}{3}x^2
%    + 2\\
%  \nabla\phi^2
%  &=
%    x^1 \csc\pi/3
%    + x^2 \cot\pi/3
%    = \dfrac{2\sqrt{3}}{3} x^1 + \dfrac{\sqrt{3}}{3} x^2
%\end{align*}
%
%Aplicado al vector $\vvv{v}=(3,2)$
%\begin{align*}
%  \nabla\phi^1
%  &=
%    -\dfrac{\sqrt{3}}{3}\left(%
%    3-\dfrac{2\sqrt{3}}{3}
%    \right)
%    + \dfrac{\sqrt{3}}{3}\, \dfrac{4\sqrt{3}}{3}
%    + 2
%    = -\sqrt{3} + \dfrac{2}{3} + \dfrac{4}{3} + 2
%    = 4 - \sqrt{3}\\
%  \nabla\phi^2
%  &=
%    \dfrac{2\sqrt{3}}{3} x^1 + \dfrac{\sqrt{3}}{3} x^2
%    = \dfrac{2\sqrt{3}}{3} \left(3-\dfrac{2\sqrt{3}}{3}\right)
%    + \dfrac{\sqrt{3}}{3}\,\dfrac{4\sqrt{3}}{3}\\
%  &=
%    2\sqrt{3}
%    - \dfrac{4}{3}
%    + \dfrac{4}{3}
%    = 2\sqrt{3}
%\end{align*}
%
%\noindent
%Resumen:
%\[
%  \text{Sistema oblicuo general }
%  \hspace{1.2em}
%  \begin{cases}
%    \begin{array}{l}
%      \nabla\phi^1 = -x^1\cot\alpha
%      + x^2\,\dfrac{1-2\cos^2\alpha}{\sin\alpha}
%      + 2\\
%      \nabla\phi^2 = x^1\csc\alpha + x^2 \cot\alpha
%    \end{array}
%  \end{cases}
%\]
%\[
%  \begin{array}{ll}
%    \text{Sistema oblicuo \qty{60}{\degree} } &
%    \begin{cases}
%      \begin{array}{l}
%        \nabla\phi^1 = -\dfrac{\sqrt{3}}{3} x^1 + \dfrac{\sqrt{3}}{3} x^2 + 2\\[2ex]
%        \nabla\phi^2 =\dfrac{2\sqrt{3}}{3} x^1 + \dfrac{\sqrt{3}}{3} x^2
%      \end{array}
%    \end{cases}\\[2ex]
%    \text{Sistema oblicuo y vector } \vvv{v} = (3,2) &
%    \begin{cases}
%      \begin{array}{l}
%        \nabla\phi^1 = 4 - \sqrt{3}\\
%        \nabla\phi^2 = 2\sqrt{3}
%      \end{array}
%    \end{cases}
%  \end{array}
%\]
%
%\noindent
%i) Cuadrado del gradiente en coordenadas oblicuas.\\
%
%Realizamos el cálculo por partes
%{\small
%\begin{align*}
%    \nabla\phi^1\nabla\phi_1
%  &=
%    \left[%
%    -x^1\cot\alpha
%    + x^2\,\dfrac{1-2\cos^2\alpha}{\sin\alpha}
%    + 2
%    \right]
%    \left(%
%    x^2\sin\alpha + 2
%    \right)\\
%  &=
%    - 2x^1\cot\alpha
%    + \left(x^2\right)^2\left(1-\cos^2\alpha\right)
%    + 2x^2\left(2\sin\alpha - \dfrac{\cos^2\alpha}{\sin\alpha}\right)
%    -x^1x^2\cos\alpha
%    + 4
%\end{align*}
%}
%{\small
%\begin{align*}
%    \nabla\phi^2\nabla\phi_2
%  &=
%    \left(x^1\csc\alpha + x^2 \cot\alpha\right)
%    \left(x^1\sin\alpha + x^2 2\sin\alpha\cos\alpha + 2 \cos\alpha\right)\\
%  &=
%    \left(x^1\right)^2
%    + 2x^1\cot\alpha
%    + 2\left(x^2\right)^2\cos^2\alpha
%    + 2x^2\dfrac{\cos^2\alpha}{\sin\alpha}
%    + 3x^1x^2\cos\alpha
%\end{align*}
%}
%
%\begin{align*}
%  |\vvv{\nabla}\phi|^2
%  &=
%    \nabla\phi^1\nabla\phi_1 + \nabla\phi^2\nabla\phi_2\\
%  &=
%    \left(x^1\right)^2
%    + \left(x^2\right)^2 (1+\cos^2\alpha)
%    + 4x^2\cos\alpha
%    + 2x^1x^2\cos\alpha
%    + 4
%\end{align*}
%
%Para nuestro sistema oblicuo


\newpage
\subsection{Ejercicio 4}
Conocido el gradiente de un campo escalar $\phi=\phi(x,y)$, en coordenadas cartesianas ortogonales
\[
  \vvv{\nabla}\phi
  = \dfrac{\partial\phi}{\partial x}\,\uvec{\i}
  + \dfrac{\partial\phi}{\partial y}\,\uvec{\j}
\]
\noindent
a) Calcula sus componentes contravariantes en un sistema de coordenadas oblicuas, cuyos
ejes forman un ángulo $\alpha$ entre sí.\\
b) Calcula sus componentes covariantes en el sistema de coordenadas oblicuas.\\

\noindent
\textbf{Solución}

\noindent
a) Componentes contravariantes en función de $\partial\phi/\partial x$ y
$\partial\phi/\partial y$.

Obtenemos los vectores $\uvec{\i}$ y $\uvec{\j}$ en función de los vectores
unitarios de la base natural $B=(\xhat{e}_1, \xhat{e}_2)$, en coordenadas oblicuas que
forman un ángulo $\alpha$, utilizando la matriz inversa de cambio de base
$\mmm{M}^{-1}$, ver \eqref{eq:R2-oblic_Minv}
\[
  \begin{pmatrix}
    \uvec{\i} & \uvec{\j}
  \end{pmatrix}
  =
  \begin{pmatrix}
    \xhat{e}_1 & \xhat{e}_2
  \end{pmatrix}
  \,
  \begin{pmatrix}
    1 & -\cot\alpha\\
    0 & \phantom{-}\csc\alpha
  \end{pmatrix}
  =
  \begin{pmatrix}
    \xhat{e}_1 & -\cot\alpha\,\xhat{e}_1 + \csc\alpha\,\xhat{e}_2
  \end{pmatrix}
\]

Así obtenemos que
\begin{align*}
  \uvec{\i} &= \xhat{e}_1\\
  \uvec{\j} &= -\cot\alpha\,\xhat{e}_1 + \csc\alpha\,\xhat{e}_2
\end{align*}

Ahora sustituimos en la expresión del gradiente en cartesianas ortogonales
\begin{align*}
  \vvv{\nabla}\phi
  &=
    \dfrac{\partial\phi}{\partial x}\,\uvec{\i}
    + \dfrac{\partial\phi}{\partial y}\,\uvec{\j}
    = \dfrac{\partial\phi}{\partial x}\,\xhat{e}_1
    + \dfrac{\partial\phi}{\partial y}
    \left(-\cot\alpha\,\xhat{e}_1+\csc\alpha\,\xhat{e}_2\right)\\
  &=
    \dfrac{\partial\phi}{\partial x}\,\xhat{e}_1
    - \cot\alpha\dfrac{\partial\phi}{\partial y}\,\xhat{e}_1
    + \csc\alpha\dfrac{\partial\phi}{\partial y}\,\xhat{e}_2\\
  &=
    \left(%
    \dfrac{\partial\phi}{\partial x}-\cot\alpha\dfrac{\partial\phi}{\partial y}
    \right)\,\xhat{e}_1
    + \csc\alpha\,\dfrac{\partial\phi}{\partial y}\,\xhat{e}_2
    = (\nabla\phi)^1\,\xhat{e}_1 + (\nabla\phi)^2\,\xhat{e}_2
\end{align*}
Así
\begin{align*}
  (\nabla\phi)^1
  &=
    \dfrac{\partial\phi}{\partial x^1}
    = \dfrac{\partial\phi}{\partial x}-\cot\alpha\dfrac{\partial\phi}{\partial y}\\
  (\nabla\phi)^2
  &=
    \dfrac{\partial\phi}{\partial x^2}
    = \csc\alpha\,\dfrac{\partial\phi}{\partial y}
\end{align*}
  
\noindent
b) Componentes covariantes en función de $\partial\phi/\partial x$ y
$\partial\phi/\partial y$.

Bajamos índices utilizando el tensor métrico
\[
  (\nabla\phi)_i = g_{ij}\,(\nabla\phi)^j
\]

\begin{align*}
  (\nabla\phi)_1
  &=
    g_{11}(\nabla\phi)^1 + g_{12}(\nabla\phi)^2
    = (\nabla\phi)^1 + \cos\alpha (\nabla\phi)^2\\
  &=
    \dfrac{\partial\phi}{\partial x}-\cot\alpha\dfrac{\partial\phi}{\partial y}
    + \cos\alpha\csc\alpha\dfrac{\partial\phi}{\partial y}
    = \dfrac{\partial\phi}{\partial x}
    - \cancelout{\cot\alpha\dfrac{\partial\phi}{\partial y}}
    + \cancelout{\cot\alpha\dfrac{\partial\phi}{\partial y}}
    =
    \dfrac{\partial\phi}{\partial x}
\end{align*}

\begin{align*}
  (\nabla\phi)_2
  &=
    g_{21}(\nabla\phi)^1 + g_{22}(\nabla\phi)^2
    = \cos\alpha (\nabla\phi)^1 + (\nabla\phi)^2\\
  &=
    \cos\alpha
    \left(%
    \dfrac{\partial\phi}{\partial x}-\cot\alpha\dfrac{\partial\phi}{\partial y}
    \right)
    + \csc\alpha\,\dfrac{\partial\phi}{\partial y}\\
  &= \cos\alpha\,\dfrac{\partial\phi}{\partial x}
    - \dfrac{\cos^2\alpha}{\sin\alpha}\,\dfrac{\partial\phi}{\partial y}
    + \dfrac{1}{\sin\alpha}\,\dfrac{\partial\phi}{\partial y}
    = \cos\alpha\,\dfrac{\partial\phi}{\partial x}
    + \sin\alpha\,\dfrac{\partial\phi}{\partial y}
\end{align*}

Así
\begin{align*}
  (\nabla\phi)_1
  &=
    \dfrac{\partial\phi}{\partial x_1}
    = \dfrac{\partial\phi}{\partial x}\\
  (\nabla\phi)_2
  &=
    \dfrac{\partial\phi}{\partial x_2}
    = \cos\alpha\,\dfrac{\partial\phi}{\partial x}
    + \sin\alpha\,\dfrac{\partial\phi}{\partial y}
\end{align*}


%%% Local Variables:
%%% mode: latex
%%% TeX-engine: luatex
%%% TeX-master: "../retazosmatematicas.tex"
%%% End:
